\documentclass{beamer}

\usetheme{metropolis} % Modern academic theme
\usecolortheme{default}
\usepackage{graphicx}
\usepackage{booktabs}
\usepackage{amsmath}
\usepackage{hyperref}
\usepackage{bm}
\usepackage{multirow}
\usepackage{wrapfig}
\usepackage{caption} % Added to define \captionof

% Highlight key terms
\newcommand{\alertterm}[1]{\alert{\textbf{#1}}}

% Title and author information
\title[Paper2Video]{
\includegraphics[height=0.6cm]{figure/logo_.png} Paper2Video: Automatic Video Generation from Scientific Papers}
\author[Z. Zhu, K. Q. Lin, M. Z. Shou]{
  Zeyu Zhu\textsuperscript{*}, Kevin Qinghong Lin\textsuperscript{*}, Mike Zheng Shou\textsuperscript{$\dag$} \\
  {\small Show Lab, National University of Singapore}
}
\date{\today}

\setbeamerfont{caption}{size=\scriptsize}

\begin{document}

%%%%%%%%%%%%%%%%%%%%%%%%%%%%%%%%%%%%
% Open Slide
%%%%%%%%%%%%%%%%%%%%%%%%%%%%%%%%%%%%
\begin{frame}
  \titlepage
  \vspace{-1cm}
  \begin{center}
    \scriptsize
    \textsuperscript{*} Equal Contribution \quad
    \textsuperscript{$\dag$} Corresponding Author \\
    \texttt{e1374425@nus.edu.sg}
  \end{center}
\end{frame}

%%%%%%%%%%%%%%%%%%%%%%%%%%%%%%%%%%%%
% Motivation
% corresponds to source: Introduction section and teaser figure
%%%%%%%%%%%%%%%%%%%%%%%%%%%%%%%%%%%%
\begin{frame}{Motivation: Need for Automated Academic Presentation Videos}
\begin{columns}
\begin{column}{0.55\textwidth}
\begin{itemize}
  \item Academic presentation videos critical for research dissemination.
  \item Manual production is highly \alertterm{labor-intensive}, requiring hours for just a few minutes.
  \item Unique challenges:
  \begin{itemize}
    \item Long-context multi-modal input from dense scientific papers.
    \item Multi-channel alignment: slides, subtitles, speech, talking-head, cursor.
    \item No existing benchmarks or tailored metrics to evaluate quality.
  \end{itemize}
  \item \alertterm{Key questions:}
  \begin{itemize}
    \item What defines a \alertterm{good academic presentation video}?
    \item How to \alertterm{automatically generate} such videos from papers?
  \end{itemize}
\end{itemize}
\end{column}
\begin{column}{0.45\textwidth}
\centering
\includegraphics[width=1\textwidth]{figure/teaser.pdf}
\end{column}
\end{columns}
\end{frame}

%%%%%%%%%%%%%%%%%%%%%%%%%%%%%%%%%%%%
% Related Work
% corresponds to source: Related Works section
%%%%%%%%%%%%%%%%%%%%%%%%%%%%%%%%%%%%
\begin{frame}{Related Work: Landscape and Challenges}
\begin{itemize}
  \item \textbf{Natural Video Generation:}
  \begin{itemize}
    \item Diffusion models generate short videos from text~\cite{sd_video,wan}.
    \item Long and multi-shot videos remain challenging~\cite{deepmind2025veo3}.
  \end{itemize}
  \item \textbf{Multi-agent frameworks for videos:}
  \begin{itemize}
    \item MovieAgent generates movies via role decomposition~\cite{wu2025automated}.
    \item PresentAgent combines slide and speech but lacks personalization and academic slide quality~\cite{shi2025presentagent}.
  \end{itemize}
  \item \textbf{AI for Research:}
  \begin{itemize}
    \item Automatic generation of scientific posters~\cite{pang2025paper2poster} and slides~\cite{zheng2025pptagent}.
    \item Few efforts on \alertterm{academic presentation video synthesis}.
  \end{itemize}
  \item \alertterm{Our work fills this gap}, focusing on long, multimodal, personalized academic video generation with comprehensive evaluation.
\end{itemize}
\end{frame}

%%%%%%%%%%%%%%%%%%%%%%%%%%%%%%%%%%%%
% Method overview slide
% corresponds to source: Method section overview
%%%%%%%%%%%%%%%%%%%%%%%%%%%%%%%%%%%%
\begin{frame}{Method Overview: The PaperTalker Agent}
\begin{figure}
    \centering
    \includegraphics[width=0.75\textwidth]{figure/method.pdf}
    \caption{Pipeline of \alertterm{PaperTalker}: Modular multi-agent approach integrating slide generation, subtitle/cursor creation, and talking-head synthesis.}
    \label{fig:method}
\end{figure}
\vspace{-0.8cm}
\begin{block}{Key Modules}
\begin{itemize}
  \item \alertterm{Slide Builder}: Generates \LaTeX~Beamer slides, refines layouts via compilation feedback.
  \item \alertterm{Subtitle \& Cursor Builder}: Produces aligned subtitles with spatially grounded cursor trajectories.
  \item \alertterm{Talker Builder}: Synthesizes personalized speech and renders talking-head videos.
  \item Supports \alertterm{slide-wise parallelization} for efficiency.
\end{itemize}
\end{block}
\end{frame}

%%%%%%%%%%%%%%%%%%%%%%%%%%%%%%%%%%%%
% Method 1: Slide Builder with layout refinement
% corresponds to source: Slide Builder subsection in method
%%%%%%%%%%%%%%%%%%%%%%%%%%%%%%%%%%%%
\begin{frame}{Method: Slide Builder and Layout Refinement}
\begin{itemize}
  \item \alertterm{Slides are generated directly} from the paper’s \LaTeX~source using Beamer for professional academic layouts.
  \item Compilation feedback (errors, warnings) used to iteratively fix code ensuring correctness.
  \item \alertterm{Tree Search Visual Choice:}
  \begin{itemize}
    \item Rule-based generation of layout variants by adjusting font and figure scales.
    \item Use vision-language model to \alertterm{select the best candidate visually}.
    \item Solves overflow and layout issues efficiently beyond parameter tuning by LLMs.
  \end{itemize}
\end{itemize}
\begin{figure}
    \centering
    \includegraphics[width=0.5\textwidth]{figure/tree_search.pdf}
    \caption{Tree Search Visual Choice for slide layout refinement.}
    \label{fig:mcts}
\end{figure}
\end{frame}

%%%%%%%%%%%%%%%%%%%%%%%%%%%%%%%%%%%%
% Method 2: Subtitle and Cursor Builder
% corresponds to source: Subtitle Builder and Cursor Builder subsections
%%%%%%%%%%%%%%%%%%%%%%%%%%%%%%%%%%%%
\begin{frame}{Method: Subtitle and Cursor Builder}
\begin{itemize}
  \item Subtitles are generated for each slide with a vision-language model.
  \item \alertterm{Visual focus prompts} highlight key content for each subtitle sentence.
  \item Cursor trajectory is \alertterm{spatially aligned} per sentence using GUI-grounding models.
  \item \alertterm{Temporal alignment} with narration timing achieved using WhisperX.
\end{itemize}
\begin{block}{Benefit}
\begin{itemize}
  \item \alertterm{Enhanced audience guidance} by highlighting the precise focus of the narration on the slides.
  \item Improves accessibility and content grounding.
\end{itemize}
\end{block}
\end{frame}

%%%%%%%%%%%%%%%%%%%%%%%%%%%%%%%%%%%%
% Method 3: Talker Builder and Parallelization
% corresponds to source: Talker Builder subsection
%%%%%%%%%%%%%%%%%%%%%%%%%%%%%%%%%%%%
\begin{frame}{Method: Talker Builder and Parallelization}
\begin{itemize}
  \item Personalize speech synthesis using author voice sample with F5-TTS at slide granularity.
  \item Talking-head video generated based on author portrait using FantasyTalking and Hallo2.
  \item \alertterm{Slide-wise parallel generation} for significant acceleration (6$\times$ speedup).
  \item You save hours of generation time while preserving video quality and engagement.
\end{itemize}
\begin{figure}
  \centering
  
\includegraphics[width=0.45\textwidth]{figure/presenter.png}
  \caption{Personalized \alertterm{talking-head} presenter output.}
  \label{fig:presenter}
\end{figure}
\end{frame}

%%%%%%%%%%%%%%%%%%%%%%%%%%%%%%%%%%%%
% Experimental Settings
% corresponds to source: Experiment section (Experimental Settings subsection)
%%%%%%%%%%%%%%%%%%%%%%%%%%%%%%%%%%%%
\begin{frame}{Experimental Settings}
\begin{itemize}
  \item \alertterm{Baselines:}
  \begin{itemize}
    \item End-to-end video generation models (e.g., Wan2.2, Veo3).
    \item Multi-agent frameworks (e.g., PresentAgent).
    \item Human-created presentations as upper bound.
  \end{itemize}
  \item \alertterm{Dataset:} Paper2Video benchmark with 101 paper-video pairs from recent AI conferences.
  \item Hardware: 8$\times$ NVIDIA RTX A6000 GPUs.
  \item Evaluation Models:
  \begin{itemize}
    \item Vision-language model: GPT-4.1.
    \item Video language model: Gemini-2.5-Flash.
  \end{itemize}
\end{itemize}
\end{frame}

%%%%%%%%%%%%%%%%%%%%%%%%%%%%%%%%%%%%
% Experimental Results: Main
% corresponds to source: Experiment Results (Main)
%%%%%%%%%%%%%%%%%%%%%%%%%%%%%%%%%%%%
\begin{frame}{Experimental Results: Main Evaluation}
\begin{table}[ht]
  \centering
  \tiny
  \caption{\textbf{Main evaluation results on Paper2Video benchmark.} Note: $\text{PaperTalker}^{*}$ is without presenter/cursor.}
  \label{tab:main_result}
  \begin{tabular}{l c c c c c}
    \toprule
    \textbf{Method} & \textbf{Sim. (Speech/Content)} & \textbf{Arena Win} & \textbf{PresentQuiz} & \textbf{IP Mem.} & \textbf{Duration (s)} \\
    \midrule
    HumanMade & \alertterm{1.00 / 5.00} & \alertterm{50\%} & 0.738 / 0.908 & - & 375 \\
    \midrule
    Wan2.2 & NA / NA & 1.1\% & 0.251 / 0.551 & 11.5\% & 4 \\
    Veo3 & 0.133 / NA & 1.2\% & 0.367 / 0.585 & 31.3\% & 8 \\
    PresentAgent & 0.045 / 1.47 & 2.0\% & 0.548 / 0.654 & 12.5\% & 430 \\
    PaperTalker$^{*}$ & \alertterm{0.646} / \alertterm{1.97} & 15.2\% & 0.835 / 0.949 & 37.5\% & 234 \\
    PaperTalker & \alertterm{0.646} / \alertterm{1.97} & \alertterm{17.0\%} & \alertterm{0.842} / \alertterm{0.951} & \alertterm{50.0\%} & 234 \\
    \bottomrule
  \end{tabular}
\end{table}
\vspace{-1.2em}
\begin{columns}
\begin{column}{0.55\textwidth}
\includegraphics[width=0.9\textwidth]{figure/human_evaluation_bar.pdf}
\end{column}
\begin{column}{0.45\textwidth}
{\scriptsize
\begin{itemize}
  \item \alertterm{PaperTalker} outperforms all baselines in both automatic and human evaluations.
  \item Presenter video boosts engagement and quality preference by approx. 10\%.
\end{itemize}
}
\end{column}
\end{columns}
\end{frame}

%%%%%%%%%%%%%%%%%%%%%%%%%%%%%%%%%%%%
% Experimental Results: Efficiency
% corresponds to source: Experimental Results: Efficiency
%%%%%%%%%%%%%%%%%%%%%%%%%%%%%%%%%%%%
\begin{frame}{Experimental Results: Efficiency and Cost}
\begin{table}[ht]
  \centering
  \small
  \caption{\textbf{Generation cost for each method.}}
  \label{table:cost}
  \setlength{\tabcolsep}{6pt}
  \begin{tabular}{lccc}
    \toprule
    Method & Token (K)$\downarrow$ & Time (min)$\downarrow$ & Cost (\$)$\downarrow$ \\
    \midrule
    Veo3~\cite{deepmind2025veo3} & NA & \textbf{0.4} & 1.667 \\
    Wan2.2~\cite{wan} & NA & 8.1 & 0.280 \\
    \midrule
    PresentAgent~\cite{shi2025presentagent} & 241 & 39.5 & 0.003 \\
    PaperTalker (w/o Talker) & \textbf{62} & 15.6 & \textbf{0.001} \\
    PaperTalker (w/o Parallel) & \textbf{62} & 287.2 & \textbf{0.001} \\
    PaperTalker & \textbf{62} & \textbf{48.1 (6$\times$ faster)} & \textbf{0.001} \\
    \bottomrule
  \end{tabular}
\end{table}
\vspace{-1em}
\begin{itemize}
  \item Parallel generation yields \alertterm{6$\times$ speedup} vs. serial.
  \item PaperTalker achieves \alertterm{lowest generation cost} among compared methods.
\end{itemize}
\end{frame}

%%%%%%%%%%%%%%%%%%%%%%%%%%%%%%%%%%%%
% Ablation experiment
% corresponds to source: Ablation Study
%%%%%%%%%%%%%%%%%%%%%%%%%%%%%%%%%%%%
\begin{frame}{Ablation Study: Impact of Cursor and Tree Search Modules}

\begin{columns}[t]
  \begin{column}{0.45\textwidth}
  \small
  \begin{itemize}
    \item \alertterm{Cursor Highlight}: Improves content localization, increasing QA accuracy from 8\% to 63\%.
    \item \alertterm{Tree Search Visual Choice}: Fixes layout overflows and improves slide design and coherence.
  \end{itemize}
  \vspace{0.3cm}
  \scriptsize
  \textbf{Table: Ablation on Cursor}
  \begin{tabular}{lc}
    \toprule
    Method & Localization Accuracy \\
    \midrule
    Without Cursor & 0.084 \\
    With Cursor & \textbf{0.633} \\
    \bottomrule
  \end{tabular}

  \vspace{0.7cm}
  \textbf{Table: Slide Quality Evaluation}
  
  \begin{tabular}{lccc}
    \toprule
    Method & Content & Design & Coherence \\
    \midrule
    HumanMade & \textbf{4.43} & \textbf{2.85} & 2.73 \\
    PPTAgent & 3.43 & 1.57 & 1.29 \\
    PaperTalker (no TS) & 4.33 & \underline{2.73} & \textbf{3.84} \\
    PaperTalker & \underline{4.34} & \textbf{2.85} & \textbf{3.84} \\
    \bottomrule
  \end{tabular}
  \end{column}
  \begin{column}{0.55\textwidth}
    \centering
    \includegraphics[width=1\textwidth]{figure/tree_search_vis.pdf}
    \captionof{figure}{Slides before (top row) and after (bottom row) Tree Search Visual Choice layout refinement.}
  \end{column}
\end{columns}
\end{frame}

%%%%%%%%%%%%%%%%%%%%%%%%%%%%%%%%%%%%
% Deficiencies and Future research
% corresponds to source: Limitations and Future Work frame and section
%%%%%%%%%%%%%%%%%%%%%%%%%%%%%%%%%%%%
\begin{frame}{Limitations and Future Research Directions}
\begin{itemize}
  \item \alertterm{Current Limitations:}
  \begin{itemize}
    \item Difficulty in parsing long scientific documents with complex equations end-to-end.
    \item Some layouts and slide designs still need refinement.
    \item Talking-head synthesis quality and naturalness can be improved further.
  \end{itemize}
  \item \alertterm{Future Directions:}
  \begin{itemize}
    \item Real-time interactive presentation video agents enabling Q\&A and dynamic tutorials.
    \item Enhanced semantic understanding of longer, complex papers.
    \item Support for multilingual presentations and accessibility features.
    \item Closed-loop evaluation with human and AI-in-the-loop feedback.
    \item Full open-sourcing for community collaboration.
  \end{itemize}
\end{itemize}
\end{frame}

%%%%%%%%%%%%%%%%%%%%%%%%%%%%%%%%%%%%
% End Slide
%%%%%%%%%%%%%%%%%%%%%%%%%%%%%%%%%%%%
\begin{frame}
  \centering
  \Huge \textbf{Thank You!}
  \vspace{0.7cm}

  
\includegraphics[width=0.15\textwidth]{figure/logo_.png}

  \vspace{0.3cm}
  \large
  \texttt{Questions and Discussion Welcome} \\
  \href{mailto:e1374425@nus.edu.sg}{e1374425@nus.edu.sg}
\end{frame}

\end{document}